\label{seminar02}
\subsection*{План семинара (90 минут)}
Работа выполняется вручную. Входные вопросы --- 5 минут; задачи С1--С2
--- 15; С3 --- 10; С4 --- 25; С5 --- 15; С6 --- 15; самопроверка --- 5.
Сначала решите условия, затем обратитесь к отдельному разделу решений.
Во всех целочисленных задачах биты нумеруются справа с нуля;
при 8-битной операции остаются 8 младших битов. Флаги задаются для
\texttt{ADD} и \texttt{SUB} архитектуры x86. Объёмы считаются без сжатия,
заголовков и выравнивания, если не сказано иное.

\subsection*{Входные вопросы}
Что задаёт смысл битов? Чем 1 KiB отличается от 1 kB?
Какая информация нужна для интерпретации \texttt{10000000}?
Почему указания «это число» недостаточно?

\subsection*{Условия обязательных задач}
\paragraph{С1. Изображение.}
Изображение имеет размер $320\times200$ пикселей.
Найдите объём RGB888 в байтах и KiB. Затем найдите объём индексов
для палитры из 16 цветов. Сколько займёт палитра, если каждый её
цвет задан тремя байтами RGB? Сравните суммарные объёмы.

\paragraph{С2. Звук.}
Запись длится 5 секунд: 48\,000 отсчётов в секунду на канал,
16 бит на отсчёт, два канала. Найдите число отсчётов одного канала,
общий объём в байтах и MB. Как изменится объём для моно?

\paragraph{С3. Текст.}
Даны кодовые точки: A --- \texttt{U+0041}, Я --- \texttt{U+042F},
знак с кодом \texttt{U+1F600}.
Для UTF-8 используйте готовые коды
\texttt{41}, \texttt{D0 AF}, \texttt{F0 9F 98 80};
для UTF-16 --- единицы \texttt{0041}, \texttt{042F}, \texttt{D83D DE00}.
Запишите байты строки из этих трёх знаков в UTF-8 и UTF-16LE без BOM.
Сколько в ней кодовых точек, единиц UTF-16 и байтов в каждой кодировке?
Что вернёт \texttt{strlen} для этой UTF-8-строки в C?

\paragraph{С4. Целые числа и флаги.}
\begin{enumerate}
    \item Укажите диапазоны 8-битного беззнакового числа и дополнительного кода.
    Представьте $-37$ в прямом, обратном, дополнительном и смещённом
    (bias $=128$) кодах.
    \item Выполните 8-битное \texttt{ADD}: $100+30$.
    Прочитайте результат как unsigned и signed. Найдите CF, OF, ZF, SF.
    \item Выполните \texttt{SUB}: $80-50$. Затем сложите $80$ с 8-битным
    кодом $-50$ командой \texttt{ADD}. Сравните биты и флаги.
    \item Самостоятельно: выполните \texttt{SUB} для $0-1$ и укажите флаги.
\end{enumerate}

\paragraph{С5. IEEE 754.}
Представьте $-6.5$ в binary32: знак, истинный порядок, поле порядка,
23 бита дроби и hex-слово. Обратно декодируйте \texttt{0xC0A00000}.
Классифицируйте слова \texttt{0x00000000}, \texttt{0x00000001},
\texttt{0x7F800000}, \texttt{0x7FC00000}, не переводя их полностью в десятичную систему.

\paragraph{С6. Маски, сдвиги, команда.}
\begin{enumerate}
    \item Для $x=\texttt{0x5A}$ проверьте бит 2, установите его,
    сбросьте бит 4 исходного $x$ и выделите биты 4--2.
    \item Для \texttt{0xA6} выполните 8-битные логический и арифметический
    сдвиги вправо на 2. Сравните signed-результат с делением $-90/4$
    с усечением к нулю.
    \item Декодируйте \texttt{0x3540} по учебному формату лекции:
    4 бита opcode, 3 dst, 3 src, 6 резервных. Opcode 3 означает
    $R_{dst}\gets R_{dst}+R_{src}$. Это x86 или условная ISA?
\end{enumerate}

\subsection*{Решения обязательных задач}
\paragraph{С1.}
Пикселей $320\cdot200=64\,000$. Объём RGB888:
\[64\,000\cdot3=192\,000\text{ байт}=187.5\text{ KiB}.\]
Для 16 цветов нужно 4 бита на индекс:
$64\,000\cdot4/8=32\,000$ байт. Палитра занимает $16\cdot3=48$ байт,
итого $32\,048$ байт. Отношение объёмов $192\,000/32\,048\approx5.99$.
Уменьшение объёма связано с ограничением до 16 выбранных цветов.

\paragraph{С2.}
В одном канале $48\,000\cdot5=240\,000$ отсчётов.
Объём: $240\,000\cdot2\cdot2=960\,000$ байт = 0.96 MB.
Первый множитель 2 --- байты на отсчёт, второй --- каналы.
Моно занимает $480\,000$ байт. Частота каждого канала остаётся 48 kHz.

\paragraph{С3.}
UTF-8: \texttt{41 D0 AF F0 9F 98 80}, 7 байт.
UTF-16LE: \texttt{41 00 2F 04 3D D8 00 DE}, 8 байт.
Кодовых точек 3, единиц UTF-16 --- 4. В этом примере видимых знаков 3.
\texttt{strlen} возвращает 7: нулевой завершающий байт в длину не входит.
Для хранения C-строки с терминатором нужно 8 байт.

\paragraph{С4.}
Диапазоны: unsigned $[0,255]$, дополнительный код $[-128,127]$.
Модуль $37=00100101_2$. Коды $-37$:
прямой \texttt{10100101}; обратный \texttt{11011010};
дополнительный \texttt{11011011}; смещённый
$-37+128=91$, то есть \texttt{01011011}.

$100+30$:
\begin{center}\texttt{01100100 + 00011110 = 10000010}.\end{center}
Unsigned-значение 130, signed-значение $130-256=-126$.
Флаги: CF=0, OF=1, ZF=0, SF=1. Переноса нет, но два положительных
операнда дали отрицательное signed-значение.

$80-50=30$: \texttt{01010000 - 00110010 = 00011110}.
Для \texttt{SUB}: CF=0 (заём не нужен), OF=0, ZF=0, SF=0.
Код $-50$ равен \texttt{11001110}, как unsigned это 206.
\texttt{ADD}: $80+206=286=256+30$, поэтому результат тоже
\texttt{00011110}, но CF=1. Остальные три флага здесь совпадают.
Одинаковые биты результата не означают одинаковые флаги разных инструкций.

$0-1$: результат \texttt{11111111}, unsigned 255, signed $-1$.
Для \texttt{SUB}: CF=1 (заём), OF=0, ZF=0, SF=1.

\paragraph{С5.}
$6.5=110.1_2=1.101_2\cdot2^2$. Поэтому $s=1$, $E=2$,
$e=129=10000001_2$, $f=\texttt{10100000000000000000000}$.
Собираем слово:
\begin{center}\texttt{1 10000001 10100000000000000000000}.\end{center}
Получаем \texttt{0xC0D00000}.
Проверка: $-(1+1/2+1/8)\cdot4=-6.5$.

Для \texttt{0xC0A00000}: $s=1$, $e=129$, дробь $0.01_2$.
Значение $-1.01_2\cdot2^2=-5$.
\texttt{0x00000000} --- $+0$;
\texttt{0x00000001} --- наименьшее положительное субнормальное $2^{-149}$;
\texttt{0x7F800000} --- $+\infty$;
\texttt{0x7FC00000} --- NaN (не единственный возможный код NaN).

\paragraph{С6.}
$x=\texttt{01011010}$, маска бита 2 --- \texttt{00000100}.
Проверка даёт 0. Установка: \texttt{01011110}=\texttt{0x5E}.
Сброс бита 4 исходного $x$: \texttt{01001010}=\texttt{0x4A}.
Биты 4--2: \texttt{110}, то есть 6.

Логический сдвиг \texttt{10100110} вправо на 2:
\texttt{00101001}=\texttt{0x29}=41.
Арифметический: \texttt{11101001}=\texttt{0xE9}=$-23$.
Он соответствует округлению $-90/4=-22.5$ вниз, тогда как целочисленное
деление C даёт $-22$ (усечение к нулю).

\texttt{0x3540 = 0011 010 101 000000}:
opcode=3, dst=2, src=5, резерв=0. Действие $R_2\gets R_2+R_5$.
Это условная ISA. Для x86 нужна другая схема декодирования.

\subsection*{Самостоятельная работа после семинара}
\begin{enumerate}
    \item Изображение $128\times128$ в RGBA8888: объём в байтах и KiB.
    \item PCM: 3 секунды, 44\,100 Гц, 16 бит, моно; найдите объём.
    \item Представьте $-18$ в 8-битном дополнительном коде и расширьте до 16 бит.
    \item Выполните \texttt{ADD} для $127+1$ и $255+1$. Сравните CF и OF.
    \item Закодируйте $0.75$ в binary32. Укажите все поля и hex-слово.
    \item Для \texttt{0x81} выполните циклический сдвиг влево на 1 в 8 битах.
    \item Дополнительно: в binary32 сравните $(a+b)+c$ и $a+(b+c)$
    при $a=2^{24}$, $b=1$, $c=-2^{24}$.
\end{enumerate}
\textbf{Краткие ответы:}\par
1) $65\,536$ байт = 64 KiB;
2) $264\,600$ байт;
3) \texttt{11101110}, \texttt{11111111 11101110};
4) \texttt{80}: CF=0, OF=1; \texttt{00}: CF=1, OF=0 (hex).\par
5) \texttt{0 01111110 10000000000000000000000}, \texttt{0x3F400000};
6) \texttt{0x03}; 7) 0 и 1 при округлении каждой операции к binary32.

\subsection*{Пять минут самопроверки}
Проверьте единицы объёма, число каналов, заданную кодировку,
ширину разрядной сетки и инструкцию при вычислении флагов.
При декодировании float сначала проверьте особые значения порядка,
а только затем применяйте формулу для нормального числа.

\subsection*{Дополнительные разобранные задачи по целым числам}
Следующие шесть задач сохраняют более подробную тренировку целочисленной
части курса. Они предназначены для самостоятельной работы и резерва
преподавателя сверх обязательного 90-минутного маршрута.
